\section{Zadanie 6}

Rzucamy \(m\) kul w sposób jednostajny i niezależny do \(n\) urn,
ponumerowanych \(0, \ldots, n-1\).
Mówimy, że w indeksie \(i\) jest \(k\)-luka,
jeśli wszystkie urny o indeksach \(i, i+1, \ldots, i+k-1\) są puste.

\begin{enumerate}[label=(\roman*)]
	\item Wyznacz wartość oczekiwaną liczby \(k\)-luk.
	\item Wyznacz nierówność Czernowa dla liczby \(k\)-luk.
\end{enumerate}

\subsection*{Rozwiązanie (i)}
Niech indykator \(X_i\) oznacza, czy w miejscu \(i\) mamy \(k\)-lukę,
czyli \(X_i = 1\) iff wszystkie urny \(i, \ldots, i+k-1\) są puste.

\( E(X_i) = P(X_i = 1) = \left(\frac{n-k}{n}\right)^{m} \)
- każda z \(m\) kul musiała trafić w inne miejsce

Mamy \((n-k+1)\) możliwych indeksów w których będzie \(k\)-luka -
\(0, 1, \ldots, n-k\)

Z liniowości wartości oczekiwanej
\[
	E\left(\sum_{i=0}^{n-k}X_i\right)
	= \sum_{i=0}^{n-k}E(X_i)
	= (n-k+1) E(X_0)
	= (n-k+1)\left(\frac{n-k}{n}\right)^{m}
\]

\subsection*{Rozwiązanie (ii)}
Przyjmijmy kilka oznaczeń.

\(X = \sum_{i=0}^{n-k}X_i\) - liczba \(k\)-luk

\(\mu = E(X) = (n-k+1)\left(\frac{n-k}{n}\right)^{m}\)

\(B = (x_0, x_1, \ldots, x_{n-1})\) - losowa tupla z
modelu \(n\) urn i \(m\) kul

\(p = (y_0, y_1, \ldots, y_{n-1})\) - niezależne zmienne o rozkładzie
Poissona z współczynnikiem \(\frac m n\)

Niech \(Z_i^{(a_0, \ldots, a_{n-1})}\) będzie liczbą \(k\)-luk na indeksach
\(\equiv i\ (\text{mod}\ k)\).

\(Z_i^B\) oraz \(Z_i^p\) będą zmiennymi losowymi, oznaczającymi liczbę
\(k\)-luk dla tupli z modelu z kulami i z modelu Poissona.

Nasze ograniczenie nie będzie szczególnie dokładne, więc możemy zignorować
podłogi i sufity i działać bezpośrednio na ułamkach.

Zgodnie z naszymi oznaczeniami:

\[
	Z_i^B = Z_i^{(x_0, \ldots, x_{n-1})}
	= \sum_{j=0}^{\frac{n}{k}} x_{i+jk}
\]
\[
	Z_i^p = Z_i^{(y_0, \ldots, y_{n-1})}
	= \sum_{j=0}^{\frac{n}{k}} y_{i+jk}
\]
\[
	E(Z_i^B) = \frac{\mu}{k}
\]

Ustalmy funkcję g - czy na konkretnym modulo (0 mod k) mamy więcej niż oczekiwane
\(k\)-luk

\[
	g((a_0, \ldots, a_{n-1})) =
	\begin{cases}
		1 & Z_0^{((a_0, \ldots, a_{n-1}))} \geqslant (1+\delta)\frac{\mu}k \\
		0 & \text{wpp}
	\end{cases}
\]

Zauważmy, że \(g\) jest monotoniczne - zmniejszając liczbę kul możemy
jedynie zwiększyć liczbę \(k\)-luk. Z twierdzenia z wykładu dla funkcji
monotonicznych (\texttt{Theorem 5.10} w P\&C) mamy:

\[
	P\left(Z_0^{B} \geqslant (1+\delta)\frac{\mu}k\right)
	= E(g(x_0, \ldots, x_{n-1}))
	\leqslant E(g(y_0, \ldots, y_{n-1})) \cdot 2
\]

\(Z_i^{p}\) jest tak skonstruowane, że te \(k\)-luki
są od siebie niezależne
- rozważamy \(k\)-luki na pozycjach modulo \(k\), więc są rozłączne.
Można traktować \(Z_i^{p}\) jako \(\frac{n}k\) niezależnych
prób Bernoulliego, a wtedy skorzystać z lematu sprzed kilku zajęć
(\texttt{Theorem 4.4} w P\&C)

\[
	E(g(y_0, \ldots, y_{n-1}))
	= P\left(Z_0^{p} \geqslant (1+\delta)\frac{\mu}k\right)
	\leqslant e^{\frac{\frac{-\mu}{k}\delta^2}{3}}
	= \exp\left(\frac{\frac{-\mu}{k}\delta^2}{3}\right)
\]

(od tego miejsca dla czytelności będę pisał \(\exp(t)\) zamiast \(e^t\))

Aby ograniczyć \(X\) (liczbę wszystkich \(k\)-luk) użyjemy upper bounda

\[
	P(X \geqslant (1+\delta)E(X))
	\leqslant P\left(\bigcup_{i=0}^{k-1}
	\left\{Z_i^B \geqslant (1+\delta)\frac{\mu}k \right\}\right)
	\leqslant^{\text{UB}}
	k P\left(Z_0^{B} \geqslant (1+\delta)\frac{\mu}k\right)
\]
\[
	\leqslant 2k\cdot P\left(Z_0^{p} \geqslant (1+\delta)\frac{\mu}k\right)
	\leqslant 2k\cdot \exp\left(\frac{\frac{-\mu}{k}\delta^2}{3}\right)
\]

Licząc lewe ograniczenie postępujemy analogicznie. Warto zwrócić uwagę,
że w \texttt{Theorem 4.5} mamy inny mianownik w wykładniku

\[
	P(X \leqslant (1-\delta)E(X))
	\leqslant P\left(\bigcup_{i=0}^{k-1}
	\left\{Z_i^B \leqslant (1-\delta)\frac{\mu}k \right\}\right)
	\leqslant^{\text{UB}}
	k P\left(Z_0^{B} \leqslant (1-\delta)\frac{\mu}k\right)
\]
\[
	\leqslant 2k\cdot P\left(Z_0^{p} \leqslant (1-\delta)\frac{\mu}k\right)
	\leqslant 2k\cdot \exp\left(\frac{\frac{-\mu}{k}\delta^2}{2}\right)
\]

{\tiny Notatka od autora:
Iterowanie się po ułamkach w tym zadaniu jest całkowicie nieformalne.
Taka jednak metoda została zaprezentowana w czasie zajęć
i taka będzie oczekiwana, gdyby podobne zadanie ewentualnie pojawiło
się w czasie kolokwium.}

\subsection*{Odpowiedź}
\[
	\text{(i)}\quad \boxed{(n-k+1)\left(\frac{n-k}{n}\right)^{m}}
\]
\[
	\text{(ii)}\quad \boxed{
		P(X \geqslant (1+\delta)E(X))
		\leqslant 2k\cdot \exp\left(\frac{\frac{-\mu}{k}\delta^2}{3}\right)
	}
\]
\[
	\text{(ii)}\quad \boxed{
		P(X \leqslant (1-\delta)E(X))
		\leqslant 2k\cdot \exp\left(\frac{\frac{-\mu}{k}\delta^2}{2}\right)
	}
\]
